% ============================================================
% NDH Citation Ontology
% Module: 11-citation-ontology.tex
% Altitude: A0–A4 (Pedagogical → Expressive → Rendering)
% Membrane: Citation-Semantics • Provenance-Safe
% ============================================================

\section{Citation Ontology}
\label{sec:citation-ontology}

The NDH Citation Ontology formalizes how external mathematical,
physical, and epistemic sources are represented within the NDH
constitutional manifold. Citations are treated as structured
geometric objects with membranes, adjacencies, and provenance
constraints. No citation may alter NDH identity, curvature,
or constitutional rules.

% ------------------------------------------------------------
% 1. Citation Manifold
% ------------------------------------------------------------

\subsection{Citation Manifold}

We define the citation manifold as a triple

\[
\mathcal{CIT} = (S, \mathcal{A}_{\mathrm{type}}, 
\mathcal{A}_{\mathrm{provenance}})
\]

where:

\begin{itemize}
    \item $S$ is the set of citation surfaces.
    \item $\mathcal{A}_{\mathrm{type}}$ is the adjacency relation
          induced by citation class (e.g., sheaf theory, category
          theory, operator algebra).
    \item $\mathcal{A}_{\mathrm{provenance}}$ is the adjacency
          relation induced by provenance constraints.
\end{itemize}

Each citation surface $s \in S$ is a structured object:

\[
s = (m, a, p)
\]

where $m$ is metadata, $a$ is adjacency class, and $p$ is provenance.

% ------------------------------------------------------------
% 2. Citation Classes
% ------------------------------------------------------------

\subsection{Citation Classes}

We define the following citation classes relevant to NDH:

\[
\mathcal{K} = 
\{\mathrm{SheafTheory}, \mathrm{CategoryTheory},
\mathrm{OperatorAlgebra}, \mathrm{ManifoldGeometry},
\mathrm{ToposQuantum}, \mathrm{ConstitutionalFunctors},
\mathrm{NAQMMPhysics}\}
\]

Each class $k \in \mathcal{K}$ induces a citation membrane:

\[
\mu_k : S \to S
\]

which restricts citation behavior according to NDH constitutional rules.

% ------------------------------------------------------------
% 3. Citation Operators
% ------------------------------------------------------------

\subsection{Citation Operators}

Citations are governed by a family of operators:

\[
\mathfrak{O}_{\mathrm{cit}} =
\{\mathcal{O}_{\mathrm{bind}},
  \mathcal{O}_{\mathrm{adj}},
  \mathcal{O}_{\mathrm{prov}},
  \mathcal{O}_{\mathrm{render}}\}
\]

with the following constraints:

\begin{itemize}
    \item $\mathcal{O}_{\mathrm{bind}}$ is \emph{non-absorptive}:
    \[
    \mathcal{O}_{\mathrm{bind}}(s) = s
    \]
    NDH does not absorb external identity.

    \item $\mathcal{O}_{\mathrm{adj}}$ preserves adjacency class:
    \[
    \mathcal{O}_{\mathrm{adj}}(s) \in \mathcal{A}_{\mathrm{type}}(s)
    \]

    \item $\mathcal{O}_{\mathrm{prov}}$ enforces provenance:
    \[
    \mathcal{O}_{\mathrm{prov}}(s) = p
    \]

    \item $\mathcal{O}_{\mathrm{render}}$ maps citations to rendering
          surfaces without altering NDH geometry.
\end{itemize}

% ------------------------------------------------------------
% 4. Constitutional Constraints
% ------------------------------------------------------------

\subsection{Constitutional Constraints}

Citations must satisfy:

\[
\forall s \in S,\quad
\pi_{\mathrm{NDH}}(s) = \mathrm{id}_{\mathcal{M}}
\]

where $\pi_{\mathrm{NDH}}$ is the identity fibration of the NDH
manifold. This ensures:

\begin{itemize}
    \item No citation may alter NDH identity.
    \item No citation may modify constitutional curvature.
    \item No citation may introduce external absorption.
\end{itemize}

% ------------------------------------------------------------
% 5. Citation Provenance
% ------------------------------------------------------------

\subsection{Citation Provenance}

Each citation surface carries provenance:

\[
p = (\mathrm{author}, \mathrm{year}, \mathrm{source}, \mathrm{access})
\]

Provenance is a first-class constitutional object and must be preserved
under all citation operators.

% ------------------------------------------------------------
% 6. Citation Stack
% ------------------------------------------------------------

\subsection{Citation Stack}

The citation ontology forms a stack:

\[
\mathfrak{CIT} =
(\mathcal{CIT}_0, \mathcal{CIT}_1, \mathcal{CIT}_2)
\]

where:

\begin{itemize}
    \item $\mathcal{CIT}_0$ governs metadata.
    \item $\mathcal{CIT}_1$ governs adjacency.
    \item $\mathcal{CIT}_2$ governs provenance.
\end{itemize}

This stack integrates with the Octo Constitution via the Pedagogical,
Expressive, and Rendering membranes.

% ------------------------------------------------------------
% End of Module
% ------------------------------------------------------------
